% ══════════════════════════════════════════════════════════════════════════════
% Lecture 18: Shortest Paths — Bellman-Ford and Dijkstra
% ══════════════════════════════════════════════════════════════════════════════

% ══════════════════════════════════════════════════════════════════════════════
\section*{The Shortest Path Problem}
\addcontentsline{toc}{subsection}{Shortest Path Problem}
% ══════════════════════════════════════════════════════════════════════════════

% ─────────────────────────────────────────────────────────────────────────────
\subsection{Problem Definition}

\begin{definition}{Shortest Path Problem}{shortestpath}
\textbf{Input:} A directed graph $G = (V, E)$ with edge weights $w: E \to \mathbb{R}$ 
that assigns a weight $w(u,v)$ to each edge $(u,v) \in E$.

The \textbf{weight of a path} $p = \langle v_0, v_1, \ldots, v_k \rangle$ is:
\[
w(p) = \sum_{i=1}^{k} w(v_{i-1}, v_i)
\]

The \textbf{shortest-path weight} from $u$ to $v$ is:
\[
\delta(u,v) = \begin{cases}
\min\{w(p) : u \leadsto v\} & \text{if a path from } u \text{ to } v \text{ exists} \\
\infty & \text{otherwise}
\end{cases}
\]

A \textbf{shortest path} from $u$ to $v$ is any path $p$ with weight $w(p) = \delta(u,v)$.
\end{definition}

% ─────────────────────────────────────────────────────────────────────────────
\subsection{Problem Variants}

There are several variants of the shortest path problem:

\begin{itemize}[noitemsep]
  \item \textbf{Single-source:} Find shortest paths from a source vertex $s$ to every other vertex.
  \item \textbf{Single-destination:} Find shortest paths to a given destination vertex from all vertices. 
        Can be solved by single-source by reversing edge directions.
  \item \textbf{Single-pair:} Find the shortest path from $u$ to $v$. No algorithm known that is 
        better in the worst case than solving single-source.
  \item \textbf{All-pairs:} Find shortest paths from $u$ to $v$ for all pairs of vertices. Can be 
        solved by running single-source for each vertex, but better algorithms exist.
\end{itemize}

We will focus on the \textbf{single-source shortest path problem}.

% ─────────────────────────────────────────────────────────────────────────────
\subsection{Negative Weights and Cycles}

We allow negative edge weights in the graph.

\begin{remark}{Negative-Weight Cycles}{negcycle}
If there is a negative-weight cycle reachable from the source, then shortest paths are 
not well-defined: we can traverse the cycle repeatedly to get paths of arbitrarily negative weight.

Some algorithms (like Dijkstra's) only work when all edge weights are non-negative. 
Others (like Bellman-Ford) can handle negative weights and detect negative cycles.
\end{remark}

\textbf{Application:} Currency exchange arbitrage. Model currencies as vertices and exchange 
rates as edge weights (using logarithms to convert products to sums). A negative cycle indicates 
an arbitrage opportunity.

% ─────────────────────────────────────────────────────────────────────────────
\subsection{Shortest Path Representation}

For each vertex $v \in V$, we maintain:
\begin{itemize}[noitemsep]
  \item $v.d$: the current \textbf{shortest-path estimate} (upper bound on $\delta(s,v)$)
  \item $v.\pi$: the \textbf{predecessor} of $v$ in the shortest path from $s$
\end{itemize}

Initially, $v.d = \infty$ for all $v \neq s$, and $s.d = 0$. The predecessor $v.\pi = \text{NIL}$ 
for all vertices.

% ══════════════════════════════════════════════════════════════════════════════
\section*{Bellman-Ford Algorithm}
\addcontentsline{toc}{subsection}{Bellman-Ford Algorithm}
% ══════════════════════════════════════════════════════════════════════════════

% ─────────────────────────────────────────────────────────────────────────────
\subsection{The Relax Operation}

The key operation in shortest path algorithms is \textbf{relaxation}.

\begin{algorithm}[H]
\caption{Relax Operation}
\begin{algorithmic}[1]
\Procedure{Relax}{$u, v, w$}
  \If{$v.d > u.d + w(u,v)$}
    \State $v.d \gets u.d + w(u,v)$ \AlgComment{Improve estimate}
    \State $v.\pi \gets u$ \AlgComment{Update predecessor}
  \EndIf
\EndProcedure
\end{algorithmic}
\end{algorithm}

\textbf{Intuition:} Can we improve the shortest path estimate for $v$ by going through $u$ 
and taking edge $(u,v)$? If $u.d + w(u,v) < v.d$, then we have found a better path to $v$.

% ─────────────────────────────────────────────────────────────────────────────
\subsection{Algorithm Description}

The Bellman-Ford algorithm works for graphs with negative weights and can detect negative cycles.

\begin{algorithm}[H]
\caption{Bellman-Ford Algorithm}
\begin{algorithmic}[1]
\Procedure{Bellman-Ford}{$G, w, s$}
  \State Initialize: $s.d \gets 0$, $v.d \gets \infty$ for all $v \neq s$, $v.\pi \gets \text{NIL}$ for all $v$
  \For{$i = 1$ to $|V| - 1$}
    \For{each edge $(u,v) \in E$}
      \State \Call{Relax}{$u, v, w$}
    \EndFor
  \EndFor
  \For{each edge $(u,v) \in E$} \AlgComment{Check for negative cycles}
    \If{$v.d > u.d + w(u,v)$}
      \State \Return \textbf{false} \AlgComment{Negative cycle exists}
    \EndIf
  \EndFor
  \State \Return \textbf{true}
\EndProcedure
\end{algorithmic}
\end{algorithm}

\textbf{Idea:} Iteratively relax all edges. After $i$ iterations, the algorithm has found 
shortest paths with at most $i$ edges. Since any simple shortest path has at most $|V|-1$ edges, 
$|V|-1$ iterations suffice.

% ─────────────────────────────────────────────────────────────────────────────
\subsection{Correctness}

\begin{theorem}{Bellman-Ford Correctness}{bfcorrect}
If the graph $G = (V, E)$ contains no negative-weight cycles reachable from the source $s$, 
then after $|V|-1$ iterations of relaxing all edges, we have $v.d = \delta(s,v)$ for all 
vertices $v \in V$.
\end{theorem}

\begin{proof}[Proof Sketch]
Consider any vertex $v$ and let $p = \langle v_0, v_1, \ldots, v_k \rangle$ be a shortest 
path from $s = v_0$ to $v = v_k$. Since there are no negative cycles, we can assume $p$ is 
simple, so $k \le |V|-1$.

Initially, $v_0.d = 0 = \delta(s, v_0)$. After $i$ iterations, we have $v_i.d = \delta(s, v_i)$ 
for all $i$ (by induction). After relaxing edge $(v_{i-1}, v_i)$ in iteration $i$, we have 
$v_i.d = v_{i-1}.d + w(v_{i-1}, v_i) = \delta(s, v_i)$. After $k \le |V|-1$ iterations, 
$v_k.d = \delta(s, v_k)$.
\end{proof}

% ─────────────────────────────────────────────────────────────────────────────
\subsection{Detecting Negative Cycles}

\begin{theorem}{Negative Cycle Detection}{negcycledetect}
There is a negative-weight cycle reachable from the source $s$ if and only if after running 
$|V|-1$ iterations of Bellman-Ford, at least one edge $(u,v)$ can still be relaxed 
(i.e., $v.d > u.d + w(u,v)$).
\end{theorem}

\begin{proof}
$(\Rightarrow)$ If a negative cycle is reachable, then shortest paths are not well-defined, 
and some $v.d$ values will continue to decrease.

$(\Leftarrow)$ If no negative cycle exists, the correctness theorem guarantees $v.d = \delta(s,v)$ 
after $|V|-1$ iterations. By the triangle inequality, $\delta(s,v) \le \delta(s,u) + w(u,v)$ 
for all edges, so $v.d \le u.d + w(u,v)$ and no edge can be relaxed.
\end{proof}

% ─────────────────────────────────────────────────────────────────────────────
\subsection{Complexity Analysis}

\begin{complexity}{Bellman-Ford Algorithm}{bellmanford}
The Bellman-Ford algorithm runs in $O(|V| \cdot |E|)$ time:
\begin{itemize}[noitemsep]
  \item Initialization: $O(|V|)$
  \item Main loop: $|V|-1$ iterations, each relaxing all $|E|$ edges: $O(|V| \cdot |E|)$
  \item Negative cycle check: $O(|E|)$
\end{itemize}
Total: $O(|V| \cdot |E|)$.
\end{complexity}

% ══════════════════════════════════════════════════════════════════════════════
\section*{Dijkstra's Algorithm}
\addcontentsline{toc}{subsection}{Dijkstra's Algorithm}
% ══════════════════════════════════════════════════════════════════════════════

% ─────────────────────────────────────────────────────────────────────────────
\subsection{Algorithm for Non-Negative Weights}

Dijkstra's algorithm is a faster algorithm for the single-source shortest path problem when 
all edge weights are non-negative.

\textbf{Idea:} Similar to Prim's algorithm for MST. Maintain a set $S$ of vertices whose 
shortest-path distances from the source are already known. Greedily grow $S$: at each step, 
add the vertex $v \notin S$ with minimum shortest-path estimate $v.d$.

\begin{algorithm}[H]
\caption{Dijkstra's Algorithm}
\begin{algorithmic}[1]
\Procedure{Dijkstra}{$G, w, s$}
  \State Initialize: $s.d \gets 0$, $v.d \gets \infty$ for all $v \neq s$, $v.\pi \gets \text{NIL}$ for all $v$
  \State $S \gets \emptyset$
  \State $Q \gets V$ \AlgComment{Min-priority queue keyed by $v.d$}
  \While{$Q \neq \emptyset$}
    \State $u \gets$ \Call{Extract-Min}{$Q$}
    \State $S \gets S \cup \{u\}$
    \For{each vertex $v \in \text{Adj}[u]$}
      \State \Call{Relax}{$u, v, w$}
      \If{$v.d$ decreased}
        \State \Call{Decrease-Key}{$Q, v, v.d$}
      \EndIf
    \EndFor
  \EndWhile
\EndProcedure
\end{algorithmic}
\end{algorithm}

% ─────────────────────────────────────────────────────────────────────────────
\subsection{Correctness}

\begin{theorem}{Dijkstra's Correctness}{dijkstracorrect}
If all edge weights are non-negative, Dijkstra's algorithm correctly computes shortest paths 
from the source $s$ to all vertices.
\end{theorem}

\begin{proof}[Proof Sketch]
\textbf{Loop invariant:} At the start of each iteration of the while loop, for all vertices 
$v \in S$, we have $v.d = \delta(s,v)$.

\textbf{Maintenance:} Let $u$ be the vertex added to $S$. We claim $u.d = \delta(s,u)$. 
Suppose not, and let $y$ be the first vertex along a shorter path to $u$ that is not in $S$. 
Let $x$ be $y$'s predecessor in $S$. Since we chose $u$ and not $y$, we have $u.d \le y.d$. 
But the path $s \leadsto x \to y$ is a prefix of a shortest path to $u$, so 
$\delta(s,y) \le \delta(s,u)$. Since weights are non-negative, 
$y.d = \delta(s,y) \le \delta(s,u) \le u.d \le y.d$, so $u.d = \delta(s,u)$.
\end{proof}

\begin{remark}{Why Non-Negative Weights?}{nonneg}
Dijkstra's algorithm fails with negative weights because once a vertex is added to $S$, we 
assume its shortest path is final. A negative edge could later provide a shorter path, 
violating this assumption.
\end{remark}

% ─────────────────────────────────────────────────────────────────────────────
\subsection{Implementation and Complexity}

Like Prim's algorithm, Dijkstra's performance depends on the priority queue implementation.

\begin{complexity}{Dijkstra's Algorithm}{dijkstra}
With a binary min-heap, Dijkstra's algorithm runs in $O(|E| \log |V|)$ time:
\begin{itemize}[noitemsep]
  \item Initialization: $O(|V|)$
  \item While loop: $|V|$ Extract-Min operations: $O(|V| \log |V|)$
  \item At most $|E|$ Decrease-Key operations: $O(|E| \log |V|)$
\end{itemize}
Total: $O(|E| \log |V|)$.

With a Fibonacci heap, the time improves to $O(|V| \log |V| + |E|)$.
\end{complexity}

\textbf{Comparison with Bellman-Ford:}
\begin{itemize}[noitemsep]
  \item Bellman-Ford: $O(|V| \cdot |E|)$, handles negative weights, detects negative cycles
  \item Dijkstra: $O(|E| \log |V|)$, requires non-negative weights only
\end{itemize}

% ─────────────────────────────────────────────────────────────────────────────
\subsection{Lecture Summary}

\begin{tcolorbox}[colback=lightblue, colframe=keyblue, fonttitle=\bfseries,
                  title={Key Points}]
\begin{itemize}[noitemsep]
  \item \textbf{Shortest Path Problem:} Find minimum-weight paths from a source vertex to all 
        other vertices in a weighted directed graph.
  \item \textbf{Relaxation:} Key operation that improves shortest-path estimates by checking 
        if going through an intermediate vertex yields a shorter path.
  \item \textbf{Bellman-Ford Algorithm:} Handles negative edge weights and detects negative cycles. 
        Time: $O(|V| \cdot |E|)$.
  \item \textbf{Correctness:} After $|V|-1$ iterations of relaxing all edges, shortest paths 
        are found (if no negative cycles exist).
  \item \textbf{Negative Cycle Detection:} If any edge can still be relaxed after $|V|-1$ 
        iterations, a negative cycle exists.
  \item \textbf{Dijkstra's Algorithm:} Greedy approach for non-negative weights. Uses priority 
        queue like Prim's MST algorithm. Time: $O(|E| \log |V|)$ with binary heap.
  \item \textbf{Correctness:} Maintains invariant that once a vertex is processed, its 
        shortest-path distance is final. Requires non-negative weights.
  \item Choose Bellman-Ford for negative weights or cycle detection; choose Dijkstra for 
        better performance when all weights are non-negative.
\end{itemize}
\end{tcolorbox}